\documentclass[11pt,a4paper]{article}
\usepackage{amsmath,amssymb,amsthm}
\usepackage{geometry}
\geometry{margin=1in}
\usepackage{hyperref}
\usepackage{enumitem}

\newtheorem{problem}{Problem}[section]

\title{\bfseries A Harmonic–Categorical Approach to the Riemann Zeta Function\\
\large Research Programme}
\author{}
\date{}

\begin{document}
\maketitle

\begin{abstract}
We outline a programme that synthesises two recent structures: (i) the \emph{greedy harmonic tree} and the associated Harmonic Sine Transform (HST), which yields a transfer operator whose Fredholm determinant is the Riemann zeta function up to an entire non‑vanishing factor; and (ii) the \emph{Quantum‑Egyptian functional equation} for Dirichlet series arising from integral modular tensor categories. The programme identifies a series of well‑defined, concrete problems whose solution would unify these perspectives and, ultimately, provide a self‑adjoint operator encoding the non‑trivial zeros of \(\zeta(s)\) as its spectrum. The Riemann Hypothesis serves as a long‑term orienting goal; the problems themselves are of independent mathematical interest.
\end{abstract}

\section{Completed Foundations}
The following results are now established and form the basis for the programme:
\begin{enumerate}[label=(\arabic*)]
    \item \textbf{Greedy harmonic tree and Haar basis.}
    The greedy expansion \(\theta/\pi = \sum_{n=0}^\infty \delta_n(\theta)/(n+2)\) defines a binary tree of intervals in \([0,\pi]\). From this tree we obtain an orthonormal Haar basis \(\mathcal{B}\) of \(L^2[0,\pi]\). The Harmonic Sine Transform \(\mathcal{H}:f\mapsto (\langle f,\psi\rangle)_{\psi\in\mathcal{B}}\) is an isometric isomorphism onto \(\ell^2(\mathcal{B})\).

    \item \textbf{Transfer operator and meromorphic continuation.}
    The map of normalised remainders \(T\) (eq.~2.2) is uniformly expanding and admits a countable Markov partition. The associated transfer operator \(\mathcal{L}_s\) on a Hardy space \(\mathcal{H}\) is trace‑class for \(\operatorname{Re}s>1\). The greedy harmonic zeta function \(\Phi(\theta,s)=\sum \delta_n(\theta)(n+2)^{-s}\) continues meromorphically to \(\mathbb{C}\) via the resolvent \((I-\mathcal{L}_s)^{-1}\); its only poles (apart from \(s=1\)) lie at the non‑trivial zeros of \(\zeta(s)\).

    \item \textbf{Fredholm determinant identity.}
    The Ruelle zeta function of \(T\) is computed as
    \[
    \det\nolimits_{(1)}(I-\mathcal{L}_s)=\frac{\zeta(s)}{\zeta(2s)}\,e^{P(s)},
    \]
    where \(P(s)\) is entire and non‑vanishing on the critical strip. Hence the zeros of \(\zeta(s)\) on the critical line are exactly the parameters for which \(1\) is an eigenvalue of \(\mathcal{L}_s\).

    \item \textbf{Explicit formula.}
    Smoothed sums of the HST coefficients satisfy a Guinand–Weil explicit formula involving the non‑trivial zeros of \(\zeta(s)\).

    \item \textbf{Quantum‑Egyptian functional equation.}
    For any integral modular tensor category \(\mathcal{C}\), the Egyptian denominators \(m_i=D^2/d_i^2\) and Galois signs \(\chi(m_i)=\pm1\) give a Dirichlet series \(\Psi_{\mathcal{C},\chi}(s)\) that satisfies a Riemann‑type functional equation (Geere, 2026). The proof uses the lattice realisation of integral MTCs and Poisson summation.
\end{enumerate}

\section{Open Problems}

The programme is divided into three phases. Each problem is stated precisely and is intended to be tractable with current techniques.

\subsection*{Phase I: Analytic Consolidation}

\begin{problem}[Rigorous determinant identification]
Provide a complete, self‑contained proof that the Fredholm determinant of \(\mathcal{L}_s\) equals \(\zeta(s)/\zeta(2s)\) times an entire non‑vanishing factor. This requires a detailed comparison of the Markov partition of \(T\) with that of the Gauss map, and the evaluation of the corresponding Ruelle zeta function along the lines of Mayer and Efrat. 
\end{problem}

\begin{problem}[Explicit residues of \(\Phi(\theta,s)\)]
Compute the residues \(R(\theta,\rho)=\operatorname{Res}_{s=\rho}\Phi(\theta,s)\) for a non‑trivial zero \(\rho\). Express them in terms of the eigenfunctions \(\psi_\rho\) of \(\mathcal{L}_\rho\) and relate the constants \(c_f(\rho)=\int f(\theta)R(\theta,\rho)\,d\theta\) to classical arithmetic data (e.g. the von Mangoldt function). 
\end{problem}

\begin{problem}[Trace‑class extension]
Prove that \(\mathcal{L}_s\) remains trace‑class (or that \(\mathcal{L}_s^2\) is trace‑class) for all \(s\) with \(\operatorname{Re}s>1/2\). Alternatively, show that a suitable regularisation yields a determinant defined on the whole plane without additional poles.
\end{problem}

\begin{problem}[Functional equation for the HST Dirichlet series]
Derive a functional equation for the HST Dirichlet series \(D_f(s)\) directly from the symmetries of the greedy tree, without relying solely on the meromorphic continuation of \(\Phi\). Can one obtain an analogue of the Quantum‑Egyptian functional equation using a theta‑function constructed from the wavelet coefficients?
\end{problem}

\subsection*{Phase II: Categorical Synthesis}

\begin{problem}[Rigorous limit \(\Psi_k \to \zeta\)]
Consider the tower of modular tensor categories \(\mathcal{C}_k = \mathrm{SU}(2)_k\). Prove that the normalised Dirichlet series \(\Psi_k(s)/\Psi_k(1/2)\) converge (in the sense of compact convergence on compact subsets) to \(\zeta(2s)/\zeta(1)\) as \(k\to\infty\). This requires analysing the distribution of the Galois signs in the trivial character case.
\end{problem}

\begin{problem}[Greedy–MTC correspondence]
Investigate whether the greedy harmonic tree can be realised as the infinite‑level limit of the fusion graph of \(\mathrm{SU}(2)_k\). More precisely, find a bijection between the nodes of the greedy tree and the simple objects of the tower, such that the harmonic fractions \(\pi/(n+2)\) correspond to scaled Egyptian denominators. If such a correspondence exists, the transfer operator \(\mathcal{L}_s\) should appear as a limit of the MTC transfer operators (resolvents) that produce \(\Psi_k(s)\).
\end{problem}

\begin{problem}[Universal limit of categorical Dirichlet series]
Generalise Problem~5 to other towers of integral MTCs (e.g. those arising from quantum groups at roots of unity). Determine whether the set of all such limits contains the Riemann zeta function as a universal attractor, and identify the necessary conditions on the sequence of categories.
\end{problem}

\subsection*{Phase III: Operator Theory and Spectral Realisation}

\begin{problem}[Spectral flow of \(\mathcal{L}_{1/2+it}\)]
Study the eigenvalues \(\lambda_j(t)\) of \(\mathcal{L}_{1/2+it}\) as functions of \(t\in\mathbb{R}\). Prove that each crossing \(\lambda_j(t)=1\) occurs at a zero of \(\zeta(1/2+it)\) and that the crossings are simple. Analyse the monotonicity of the eigenvalues and determine whether the Riemann Hypothesis is equivalent to the reality of the crossings.
\end{problem}

\begin{problem}[Unitary group and its generator]
Define the unitary operator \(U_t = \mathcal{L}_{1/2+it} \mathcal{L}_{1/2-it}^{-1}\) (where defined). Show that \(U_t\) is a strongly continuous one‑parameter unitary group. By Stone’s theorem, \(U_t = e^{i t H}\) for a (possibly unbounded) self‑adjoint operator \(H\). Prove that the spectrum of \(H\) coincides with the set of \(\gamma\) such that \(\zeta(1/2+i\gamma)=0\). This would be a concrete Hilbert–Pólya realisation.
\end{problem}

\begin{problem}[Finite‑dimensional approximations]
For each \(k\), construct a self‑adjoint matrix \(H_k\) (e.g. from the MTC data of \(\mathrm{SU}(2)_k\)) whose eigenvalues approximate the imaginary parts of the zeros of the corresponding Dirichlet polynomial \(\Psi_k(1/2+it)\). Prove that these matrices converge in the sense of spectral measures to a limiting self‑adjoint operator whose spectrum is the set of Riemann zeros.
\end{problem}

\begin{problem}[Dirac operator approach]
Explore the possibility of building a Dirac‑type operator directly on the greedy harmonic tree, similar to the Connes–Consani approach for the Riemann zeta function. The tree structure provides a natural Laplacian and a spin structure; the eigenvalues of the Dirac operator may encode the zeros.
\end{problem}

\section{Epilogue: The Riemann Hypothesis as a North Star}
None of the problems listed above assumes the Riemann Hypothesis. They are concrete, well‑posed mathematical questions whose solutions would deepen our understanding of the zeta function and its zeros. If the programme is fully executed, it will produce a self‑adjoint operator whose spectrum is the set of non‑trivial zeros. The reality of the eigenvalues would then be a theorem, establishing the Riemann Hypothesis as a consequence of the constructed spectral theory. Until that point, the Riemann Hypothesis serves as the ultimate motivation and a guiding principle, not a prerequisite.

\begin{thebibliography}{9}
\bibitem{Geere2026a}
V.~Geere, \emph{A Purely Geometric Reconstruction of the Sine Function}, \url{https://victorgeere.co.za/maths/sine-harmonic.html}.
\bibitem{Geere2026b}
V.~Geere, \emph{On the Correlation Structure of a Greedy Harmonic Decomposition of the Sine Function}, \url{https://victorgeere.co.za/maths/greedy-harmonic-decomposition.html}.
\bibitem{Geere2026c}
V.~Geere, \emph{The Quantum‑Egyptian Functional Equation}, v8, \url{https://victorgeere.co.za/maths/quantum-egyptian-functional-equation-v8.html}.
\bibitem{Mayer}
D.~H.~Mayer, \emph{The thermodynamic formalism approach to Selberg's zeta function for \(\mathrm{PSL}(2,\mathbb{Z})\)}, Bull. Amer. Math. Soc. (N.S.) 25 (1991), 55--60.
\bibitem{Efrat}
I.~Efrat, \emph{Dynamics of the continued fraction map and the spectral theory of \(\mathrm{SL}(2,\mathbb{Z})\)}, Invent. Math. 114 (1993), 207--218.
\bibitem{Weil52}
A.~Weil, \emph{Sur les formules explicites de la th\'eorie des nombres}, Comm. S\'em. Math. Univ. Lund (1952), 252--265.
\bibitem{Polya}
G.~P\'olya, \emph{On the zeros of certain trigonometric integrals}, J. London Math. Soc.~1 (1926), 98--99.
\bibitem{ConnesConsani}
A.~Connes, C.~Consani, \emph{Schemes over \(\mathbb{F}_1\) and zeta functions}, Compos. Math. 146 (2010), 1383--1415.
\end{thebibliography}

\end{document}